% pdflatex-tate-float.tex -- tate 浮遊体（figure/table）フィクスチャ
% uplatex+tate を ground truth とした実測比較の結論:
%   * [t] 浮遊体は本文カラムの右側カラムに置かれ，内容は縦組みとして
%     組まれる（キャプション文・和文は正立縦積み，欧文文字は横倒し縦積み）．
%     pdftex回転出荷も uplatex もカラム位置が一致（約1.5pt以内）．
%   * \includegraphics 画像も双方とも「回転されて」出る（uplatex: cm
%     0 -60 40 0; pdftex: shipout回転内）→ 正立ボックス化は行わない
%     （JLREQ的にも縦組本の図版は縦枠に従う；uplatexをGTとする限り現状維持）．
% 検証内容:
%   FLOAT-FNUM    : \figurename=図, \tablename=表 が空でなく期待文字
%   FLOAT-CAPTYPE : figure 環境内で \@captype=figure
%   FLOAT-COUNT   : \caption で c@figure=1, c@table=1 が立つ
%   FLOAT-JAREG   : キャプション相当サイズ(\smallの箱)でも \ptj@zwreg が
%                   正しく追従（旧連結バグの回帰ガード）
%   コンパイル成立自体が [t]/[b] フロート + \caption + tabular のゲート
\documentclass[pdflatex,tate,paper=a5]{jlreq}
\pagestyle{empty}
\begin{document}
\makeatletter
% --- FLOAT-FNUM: 図表名（tatechuyoko* はパススルー：図＊/表＊の＊部分のみ）
\protected@edef\jlreq@testa{\figurename}%
\def\jlreq@testb{図}%
\ifx\jlreq@testa\jlreq@testb
  \typeout{TATE-PASS: FLOAT-FNUM-fig (\jlreq@testa)}%
\else
  \typeout{TATE-FAIL: FLOAT-FNUM-fig (\jlreq@testa)}%
\fi
\protected@edef\jlreq@testa{\tablename}%
\def\jlreq@testb{表}%
\ifx\jlreq@testa\jlreq@testb
  \typeout{TATE-PASS: FLOAT-FNUM-tab (\jlreq@testa)}%
\else
  \typeout{TATE-FAIL: FLOAT-FNUM-tab (\jlreq@testa)}%
\fi
% --- FLOAT-JAREG: \small 箱での実寸追従（\small = .9\jlreq@fontsize）
\setbox\@tempboxa=\vbox{\reset@font\small
  \xdef\jlreq@testjareg{\the\ptj@zwreg}X}%
\ifdim\jlreq@testjareg=\dimexpr.9\jlreq@fontsize\relax
  \typeout{TATE-PASS: FLOAT-JAREG (\jlreq@testjareg)}%
\else
  \typeout{TATE-FAIL: FLOAT-JAREG (\jlreq@testjareg, expected \the\dimexpr.9\jlreq@fontsize\relax)}%
\fi
\makeatother
縦組み本文の中に配置される浮遊体の前。
\begin{figure}[t]
\centering
\rule{30pt}{20pt}
\caption{「浮遊体」のキャプション、禁則文字を含む長めのもの。}
\makeatletter
\protected@edef\jlreq@testa{\@captype}%
\def\jlreq@testb{figure}%
\ifx\jlreq@testa\jlreq@testb
  \typeout{TATE-PASS: FLOAT-CAPTYPE (\jlreq@testa)}%
\else
  \typeout{TATE-FAIL: FLOAT-CAPTYPE (\jlreq@testa)}%
\fi
\makeatother
\end{figure}
浮遊体の後の本文。表組みの浮遊体も試す。
\begin{table}[b]
\centering
\caption{表組みテスト}
\begin{tabular}{ll}
項目一 & 値 \\
項目二 & 値 \\
\end{tabular}
\end{table}
\makeatletter
\ifnum\c@figure=1\relax
  \typeout{TATE-PASS: FLOAT-COUNT-fig (\the\c@figure)}%
\else
  \typeout{TATE-FAIL: FLOAT-COUNT-fig (\the\c@figure, expected 1)}%
\fi
\ifnum\c@table=1\relax
  \typeout{TATE-PASS: FLOAT-COUNT-tab (\the\c@table)}%
\else
  \typeout{TATE-FAIL: FLOAT-COUNT-tab (\the\c@table, expected 1)}%
\fi
\makeatother
\end{document}
